% Supplementary tables and figures (not counted in the 6-page limit: the guidelines exclude appendices).
% Moved here from the main sections on 16 Nov (and 21 Nov, to fit 6 pages); every one is referenced from the
% main text.
\section{Supplementary results}
\label{app:supplementary}
These tables and figures support the main text; each is referred to from the section that discusses it.

\begin{figure}[!htbp]
\centering
\includegraphics[width=\columnwidth]{delivery_time_by_prebook}
\caption{Pre-booked orders follow the slot the customer picked, not the delivery (all days, before
cleaning), and hold almost every delivery over 180 minutes, so they are excluded.}
\label{fig:prebook}
\end{figure}

\begin{figure}[!htbp]
\centering
\includegraphics[width=\columnwidth]{promise_padding}
\caption{Actual minus promised delivery time (training days): most orders arrive early, and arrivals
drop sharply at +8 minutes, so promises are also compared at equal average length.}
\label{fig:promise_padding}
\end{figure}

\begin{figure}[!htbp]
\centering
\includegraphics[width=\columnwidth]{leakage_audit}
\caption{Cross-validated MAE of the gradient-boosting model when after-order times are added (training
days). The pickup time alone gives almost all of the gain, but it is not known when we predict.}
\label{fig:leakage}
\end{figure}

\begin{figure}[!htbp]
\centering
\includegraphics[width=\columnwidth]{final_test_gap}
\caption{Late-rate gap (our q80 promises minus the platform's, percentage points) as both are shifted to
the same mean promised time; below 0 means we are late less often. Dotted line: the platform's actual
mean promise. Test days.}
\label{fig:promise_gap}
\end{figure}

\begin{figure}[!htbp]
\centering
\includegraphics[width=\columnwidth]{hgb_complexity_curve}
\caption{Training vs cross-validation MAE for gradient boosting as the number of leaves per tree grows
(other settings fixed: 300 trees, learning rate 0.1, min.\ leaf 5; bars: spread over the four folds).
Validation error is lowest around 16 leaves; at 511 leaves training error is 2.45 minutes but validation
error 6.63, a clear sign of overfitting. The final model (Table~\ref{tab:models}) uses 76 leaves with a
lower learning rate, larger leaves and strong L2, which hold this overfitting back.}
\label{fig:complexity}
\end{figure}

\begin{figure}[!htbp]
\centering
\includegraphics[width=\columnwidth]{day_profiles}
\caption{Median delivery time by hour for every day (left: weekdays in blue, Saturday and Sunday in black)
and for the two Mondays a week apart (right). Shaded hours have under 2\% of orders. Includes the test
days; drawn after the test run.}
\label{fig:days}
\end{figure}

\begin{figure}[!htbp]
\centering
\includegraphics[width=\columnwidth]{rider_pressure}
\caption{Orders waiting per available rider (median) by district and hour, training days. Available riders
are those who finished a delivery in the district in the last 30 minutes; grey: no orders.}
\label{fig:riders}
\end{figure}

\begin{figure}[!htbp]
\centering
\includegraphics[width=\columnwidth]{partial_dependence}
\caption{Partial dependence of the extended gradient boosting model (all extra features, fitted on the six
training days): the average prediction for \num{5000} training orders when one input is set to each value
and the rest are kept. Distance is shown up to its 99th percentile.}
\label{fig:pd}
\end{figure}

\begin{table}[t]
\centering
\caption{Unpadded vs unpadded on the test days: our estimate against the platform's promise minus its average padding (8.20 min, learned on the training days only; second row: the average padding of the same distance band, which is fairer to the platform). MAE in minutes; $\Delta$ = ours $-$ de-padded promise with a paired bootstrap 95\% interval. Segment rows use the per-band padding.}
\label{tab:depadded}
\small
\begin{tabular}{@{}lrccc@{}}
\toprule
Orders & $n$ & Ours & Promise $-$ padding & $\Delta$ \\
\midrule
All test orders & \num{139787} & 6.49 & 6.48 & +0.02 [-0.01, +0.04] \\
All, padding per distance band & \num{139787} & 6.49 & 6.36 & +0.13 [+0.11, +0.16] \\
\midrule
Under 1 km & \num{42785} & 5.61 & 5.86 & -0.25 [-0.29, -0.20] \\
1-3 km & \num{86939} & 6.67 & 6.19 & +0.48 [+0.45, +0.51] \\
Over 3 km & \num{10063} & 8.73 & 9.94 & -1.20 [-1.38, -1.04] \\
Lunch peak (11:00) & \num{23058} & 7.22 & 6.68 & +0.53 [+0.46, +0.61] \\
Sunday 23 & \num{75063} & 6.77 & 6.47 & +0.31 [+0.27, +0.35] \\
Monday 24 & \num{64724} & 6.17 & 6.23 & -0.07 [-0.11, -0.03] \\
\bottomrule
\end{tabular}
\end{table}


\begin{table}[t]
\centering
\caption{Extension (CV only): frozen gradient boosting with extra order-time features. Difference in MAE against the base features on the same orders, with paired bootstrap 95\% intervals. The test days were not used.}
\label{tab:extension}
\small
\begin{tabular}{@{}lccc@{}}
\toprule
Features & CV MAE & $\Delta$ vs base & Lunch MAE \\
\midrule
Base (final model) & 6.25 & -- & 7.04 \\
Base + rider availability & 6.20 & -0.05 [-0.05, -0.05] & 7.03 \\
Base + restaurant history & 6.01 & -0.24 [-0.25, -0.23] & 6.82 \\
Base + both & 5.97 & -0.29 [-0.29, -0.28] & 6.81 \\
\bottomrule
\end{tabular}
\end{table}


\begin{table}[t]
\centering
\caption{Promise setting at the same mean promised time (38.2 min; our q80 promises shifted by +2.7 min). Test days. Gap = ours $-$ platform, 95\% bootstrap CI.}
\label{tab:promise}
\begin{tabular}{lrrl}
\toprule
Metric & Ours & Platform & Gap \\
\midrule
Late $>$0 min & 15.3\% & 14.9\% & +0.3 pp [+0.2, +0.5] \\
Late $>$5 min & 6.1\% & 5.0\% & +1.2 pp [+1.1, +1.3] \\
Late $>$10 min & 2.4\% & 1.5\% & +0.9 pp [+0.9, +1.0] \\
Pinball loss ($q=0.8$) & 2.52 & 2.36 & -- \\
Coverage & 84.7\% & 85.1\% & -- \\
\bottomrule
\end{tabular}
\end{table}


% F7 (instructor feedback): how industry estimates delivery times, next to ours. Hand-written (Lynn, 16 Nov)
% from sources read on 16 Nov 2026; only facts stated in those sources. Goes in Related Work (related.tex).
\begin{table*}[t]
\centering
\caption{How deployed systems estimate delivery or arrival times, next to ours. Errors are not comparable across
rows: each system uses its own data, target and baseline.}
\label{tab:benchmark}
\footnotesize
\begin{tabular}{@{}p{0.13\textwidth}p{0.17\textwidth}p{0.29\textwidth}p{0.33\textwidth}@{}}
\toprule
System & Model & Inputs at prediction time & Reported result \\
\midrule
Ele.me (Alibaba), order fulfilment time~\cite{zhu2020fulfillment,alibaba2020fulfillment} &
Deep network (encoder + predictor), attention over nearby couriers &
Area, restaurant, time, order size; real-time supply--demand ratio; orders waiting at the restaurant; nearby
couriers' distance and load; restaurant history &
Online A/B test: MAE $-9.8\%$ and complaints $-19.3\%$ against the previous production model \\
Uber, DeepETA~\cite{uber2022deepeta,hu2022deepreta} &
Linear transformer correcting a routing-engine ETA (residual) &
Origin, destination, time of request, real-time traffic, request type (e.g.\ delivery drop-off) &
Lower error than the previous XGBoost model; asymmetric Huber loss so that being late can cost more than
being early \\
Ours &
Gradient boosting (+ Elastic-Net, Random Forest) &
Distance, time of day, district, prep estimate, recent load; public data only &
Test MAE 6.49 min, level with the platform's promise minus its average 8.2-min padding (6.48); slightly behind with per-distance padding (6.36); extra
real-time features: CV MAE $-0.33$ min (5.92) \\
\bottomrule
\end{tabular}
\end{table*}


\begin{table}[t]
\centering
\caption{Elastic-Net effects in minutes of delivery time per unit of each input, holding the others fixed (fitted on the six training days). Correlated inputs share their effect, so read signs and sizes, not causes. A default prep estimate sets both prep flags, so its total effect is their sum (base: $+2.0$). Time of day (sine and cosine) and districts are left out.}
\label{tab:effects}
\small
\begin{tabular}{@{}lrr@{}}
\toprule
Input & Base & Extended \\
\midrule
Distance (per km) & +6.09 & +5.76 \\
Prep estimate (per min) & +0.66 & +0.58 \\
Prep estimate is a default & -0.54 & +1.53 \\
Prep estimate missing & +2.54 & 0.00 \\
Lunch peak (11:00) & +4.80 & +3.91 \\
Dinner peak (17--19) & +3.58 & +2.73 \\
District load (per $\times$ usual) & +1.14 & +0.43 \\
Orders per available rider (per order) & -- & +0.44 \\
Restaurant's past delivery time (per min) & -- & +0.13 \\
Restaurant's past pickup time (per min) & -- & +0.17 \\
Restaurant's prep-estimate bias (per min) & -- & +0.46 \\
District's recent delivery time (per min) & -- & +0.17 \\
Orders waiting at the restaurant (per order) & -- & +0.30 \\
\bottomrule
\end{tabular}
\end{table}


\begin{table}[t]
\centering
\caption{Simple to complex models on the base and the extended features (CV only, frozen settings, the same orders). $\Delta$ with paired bootstrap 95\% intervals.}
\label{tab:ladder}
\small
\begin{tabular}{@{}lcccc@{}}
\toprule
Model & Base & Extended & $\Delta$ & $R^2$ \\
\midrule
Elastic-Net & 6.60 & 6.27 & -0.33 [-0.34, -0.32] & 0.47 $\to$ 0.52 \\
Random Forest & 6.32 & 6.00 & -0.31 [-0.32, -0.30] & 0.51 $\to$ 0.57 \\
Gradient boosting & 6.25 & 5.92 & -0.33 [-0.34, -0.32] & 0.51 $\to$ 0.57 \\
\bottomrule
\end{tabular}
\end{table}


\begin{table}[t]
\centering
\caption{CV MAE (minutes) by segment: the platform's promise minus its per-distance padding (learned on earlier days), our final model, and our best extended model (all of the above + prep bias).}
\label{tab:extension_platform}
\small
\begin{tabular}{@{}lrccc@{}}
\toprule
Orders & Share & Platform & Base & Extended \\
\midrule
All orders & 100\% & 6.24 & 6.25 & \textbf{5.92} \\
Under 1 km & 30\% & 5.82 & 5.44 & \textbf{5.14} \\
1-3 km & 63\% & 6.09 & 6.42 & \textbf{6.08} \\
Over 3 km & 7\% & 9.44 & 8.23 & \textbf{7.91} \\
Lunch peak (11:00) & 17\% & \textbf{6.66} & 7.04 & 6.75 \\
Dinner peak (17-19) & 30\% & 6.44 & 6.67 & \textbf{6.30} \\
Saturday 22 & 27\% & 6.30 & 6.44 & \textbf{6.01} \\
\bottomrule
\end{tabular}
\end{table}

